\begin{abstract}

Automatic recognition of herbal plant species is an important computer vision task with applications in biodiversity conservation, agriculture, and traditional medicine. However, accurate plant classification remains challenging due to variations in lighting conditions, viewing angles, image quality, and background complexity. This paper presents a comparative study of traditional computer vision and deep learning approaches for herbal plant classification using a dataset constructed from 21 video sequences collected under diverse environmental conditions. A total of 425 images representing 14 plant classes were extracted and preprocessed through vegetation-based filtering, blur detection, image normalization, and data augmentation. Three classification strategies were investigated: (1) a traditional machine learning pipeline based on ORB feature extraction and Bag-of-Visual-Words representation, (2) a custom Convolutional Neural Network (CNN), and (3) transfer learning using MobileNetV2 pre-trained on ImageNet. Experimental results show that the traditional ORB-BoVW-SVM approach achieved a validation accuracy of 47.92\% and a test accuracy of 36.46\%, while the custom CNN achieved 22.40\% test accuracy. In contrast, the transfer learning model significantly outperformed both approaches, achieving a test accuracy of 93.23\%. These findings demonstrate the effectiveness of transfer learning for plant recognition tasks with limited training data and highlight the advantages of leveraging pre-trained deep feature representations over handcrafted features and shallow neural networks. The proposed framework provides a reproducible and effective solution for automated herbal plant identification in real-world environments.

\end{abstract}
